\chapter{Seguimiento de Bill y comparación de controladores}

Antes de integrarlo en la celda completa, se probó el seguimiento de Bill en
\path{Actividad2_1_Basic_Follower.ttt} (Fig.~\ref{fig:follower-replay},
capturas en Fig.~\ref{fig:follower-coppeliasim-screenshots}). Bill camina con
aceleración y frenado, y el Pioneer sigue un punto detrás de la persona. Si se
apunta al centro del cuerpo, el robot tiende a acercarse demasiado; por eso la
referencia va detrás de Bill.

\begin{figure}[H]
  \centering
  \includegraphics[width=\textwidth]{follower_replay_clean.png}
  \caption{Escena de seguimiento: Bill define un objetivo móvil y R1 regula distancia, rumbo y suavidad de ruedas.}
  \label{fig:follower-replay}
\end{figure}

\begin{figure}[H]
  \centering
  \begin{subfigure}[b]{0.49\textwidth}
    \centering
    \includegraphics[width=\textwidth]{A2_T1_Follower_SS1.png}
    \caption{Vista general de la escena.}
  \end{subfigure}
  \hfill
  \begin{subfigure}[b]{0.49\textwidth}
    \centering
    \includegraphics[width=\textwidth]{A2_T1_Follower_SS2.png}
    \caption{Seguimiento sobre la trayectoria circular.}
  \end{subfigure}
  \caption{Capturas de ejecución de \texttt{Actividad2\_1\_Basic\_Follower.ttt}: Bill recorre una trayectoria circular y R1 mantiene una distancia de seguimiento estable.}
  \label{fig:follower-coppeliasim-screenshots}
\end{figure}

\section{Geometría de seguimiento}

Si \(p_B=[x_B,y_B]^T\) es la posición de Bill, \(\theta_B\) su orientación y
\(d_{des}\) la separación deseada, el punto de referencia se calcula como:

\[
p^\star =
p_B -
d_{des}
\begin{bmatrix}
\cos\theta_B\\
\sin\theta_B
\end{bmatrix},
\qquad d_{des}=\SI{0.82}{m}.
\]

El error de distancia se mide con respecto a Bill, mientras que el error de
rumbo se calcula desde el robot hacia Bill (Fig.~\ref{fig:follower-geometry}):

\[
e_d = \|p_B-p_R\|-d_{des},
\qquad
e_\psi = \operatorname{wrap}\left(\operatorname{atan2}(y_B-y_R,x_B-x_R)-\theta_R\right).
\]

\begin{figure}[H]
  \centering
  \includegraphics[width=0.92\textwidth]{a2_follower_geometry.png}
  \caption{Referencia móvil del seguidor: el robot controla un punto detrás de Bill y no el centro de la persona.}
  \label{fig:follower-geometry}
\end{figure}

\section{Modelo cinemático usado por todos los modos}

El movimiento plano del Pioneer se aproxima con el modelo unicycle:

\[
\dot{x}=v\cos\theta,\qquad
\dot{y}=v\sin\theta,\qquad
\dot{\theta}=\omega.
\]

La conversión a velocidades angulares de rueda usa el radio \(r\) y la distancia
entre ruedas \(L\):

\[
\omega_L = \frac{v-\frac{L}{2}\omega}{r},
\qquad
\omega_R = \frac{v+\frac{L}{2}\omega}{r},
\qquad
r=\SI{0.0975}{m},\quad L=\SI{0.331}{m}.
\]

\label{sec:energy-model}
Estas mismas señales se usan para calcular potencia aproximada, energía,
variación de comando y suavidad de control.

La energía de la Tabla~\ref{tab:follower-controller-metrics} se calcula como una
métrica mecánica aproximada. Se obtiene en Lua con la potencia de avance de cada
rueda, fricción de rodadura, un término viscoso de rueda y potencia de guiñada:

\[
\begin{aligned}
P_L &= \left|\frac{m}{2}a\,v_L\right|
      + \frac{F_r}{2}|v_L| + b_\omega\omega_L^2,\\
P_R &= \left|\frac{m}{2}a\,v_R\right|
      + \frac{F_r}{2}|v_R| + b_\omega\omega_R^2,\\
P_\theta &= |I_z\dot{\omega}\omega|,\qquad
E = \sum_k (P_L+P_R+P_\theta)\Delta t .
\end{aligned}
\]

Los parámetros usados fueron \(m=\SI{16.5}{kg}\), \(I_z=\SI{0.72}{kg.m^2}\),
\(F_r=\SI{2.8}{N}\) y \(b_\omega=\SI{0.006}{kg\,m^2\,s^{-1}}\) (amortiguamiento viscoso rotacional, unidades SI: N\,m\,s). En esa expresión, \(a=\dot{v}\) es la
aceleración lineal estimada entre muestras y \(\dot{\omega}\) es la aceleración
angular. Esta magnitud se usa como indicador relativo entre controladores dentro
de la misma escena; no representa una estimación eléctrica de batería.

Los términos \(\frac{m}{2}a\,v_L\) y
\(\frac{m}{2}a\,v_R\) reparten aproximadamente la potencia inercial lineal,
no como un modelo dinámico completo de rueda. La inercia real actúa sobre el
centro de masa y se transmite por los contactos; se reparte por rueda solo
para obtener una métrica relativa sensible a frenazos, inversiones de sentido y
oscilaciones de comando.

\section{P, PI y PID}

Los controladores clásicos se aplican sobre \(e_d\) y se combinan con una
realimentación de rumbo:

\[
u_P=K_p e_d,
\qquad
u_{PI}=K_p e_d + K_i\sum_{k=0}^{n} e_d(k)\Delta t,
\]

\[
u_{PID}=K_p e_d + K_i\sum_{k=0}^{n} e_d(k)\Delta t
+ K_d\frac{e_d(k)-e_d(k-1)}{\Delta t}.
\]

El término proporcional responde al error instantáneo; el integral acumula
el retardo sistemático del seguidor y lo elimina en régimen permanente,
tanto si el objetivo se mueve como si está detenido; el derivativo amortigua
cambios bruscos del error. En ejecución, todos los modos añaden una corrección de evitación
de obstáculos al ángulo de giro:
\[
\omega = \omega_{\text{controlador}} + \omega_{obs},
\]
donde \(\omega_{obs}\) es el término repulsivo de la sección de campos potenciales
(\S\ref{ch:campos-potenciales}). Todos los modos pasan por saturaciones de velocidad y aceleración.

\section{LQR discreto}

El modo LQR del seguidor regula el error local alrededor del punto \(p^\star\),
siguiendo la forma clásica de control cuadrático lineal \cite{anderson1990}. El
estado se define en el marco del robot mediante la transformación:

\[
e_{long} = \cos\theta_R\,(x_{p^*}-x_R) + \sin\theta_R\,(y_{p^*}-y_R),
\]
\[
e_{lat}  = -\sin\theta_R\,(x_{p^*}-x_R) + \cos\theta_R\,(y_{p^*}-y_R),
\qquad
e_{yaw}  = \operatorname{wrap}(\psi_{ref}-\theta_R),
\]

donde \((x_{p^*},y_{p^*})\) es el punto de referencia \(p^\star\). El vector de
estado y la acción de control son:

\[
x_L =
\begin{bmatrix}
e_{long} & e_{lat} & e_{yaw}
\end{bmatrix}^T,
\qquad
u=
\begin{bmatrix}
\Delta v & \Delta\omega
\end{bmatrix}^T.
\]

Con \(v_{ref}=\SI{0.22}{m/s}\), radio local de giro \(\rho_{ref}=\SI{1.55}{m}\),
\(\omega_{ref}=v_{ref}/\rho_{ref}\) y
\(\Delta t=\SI{0.05}{s}\), la linealización usada en
\path{compute_follower_lqr_gain.py} es:

\[
A=
\begin{bmatrix}
0 & \omega_{ref} & 0\\
-\omega_{ref} & 0 & v_{ref}\\
0 & 0 & 0
\end{bmatrix},
\qquad
B=
\begin{bmatrix}
-1 & 0\\
0 & 0\\
0 & -1
\end{bmatrix}.
\]

La segunda fila de \(B\) es \([0,\;0]\) porque en la linealización del unicycle
alrededor de la trayectoria de referencia en el marco de Frenet, el error
lateral satisface \(\dot{e}_{lat}=v_{ref}e_{yaw}-\omega_{ref}e_{long}\):
ninguna entrada \((\Delta v,\Delta\omega)\) actúa directamente sobre \(e_{lat}\).
El error lateral solo se puede reducir indirectamente, a través del canal
de error de rumbo \(e_{yaw}\), lo que el LQR compensa asignando un peso alto
a \(e_{lat}\) en la matriz \(Q\) (\(Q_{22}=80\)).

\[
A_d=I+\Delta t A,\qquad B_d=\Delta t B.
\]

El script resuelve la ecuación algebraica de Riccati discreta con:

\[
Q=\operatorname{diag}(35,80,18),
\qquad
R=\operatorname{diag}(6,3),
\]

\[
K=(B_d^T P B_d+R)^{-1}B_d^T P A_d.
\]

Los pesos proceden de una calibración empírica sobre el seguimiento de Bill, no
de una identificación formal del Pioneer. Se penalizó más el error lateral
\((q_{lat}=80)\) que el longitudinal \((q_{long}=35)\), porque el robot admite
permitirse ir algo adelantado o retrasado, pero no cruzarse lateralmente con la
referencia móvil. El coste de giro en \(R\) se dejó menor que el de avance para
que el regulador prefiriera corregir rumbo antes de cambiar mucho la velocidad.

La ley aplicada en Lua usa \(u=-Kx_L\), con saturación posterior de \(v\) y
\(\omega\).

La matriz integrada en Lua fue:

\[
K=
\begin{bmatrix}
-2.3496 & 1.2925 & 0.1203\\
0.2427 & -4.3570 & -2.6887
\end{bmatrix}.
\]

Esta parte corresponde a un LQR discreto real, de orden bajo y aplicado al
seguimiento local de Bill.

\section{NMPC por disparo directo}
\label{sec:nmpc}

El modo NMPC se implementa como control predictivo no lineal por disparo
directo. Mantiene una secuencia de \(v_k,\omega_k\), simula el modelo hacia
adelante y mejora esa secuencia con dos pasadas de búsqueda local, siguiendo la
estructura general de MPC \cite{rawlings2017}. La optimización se resuelve sobre
la secuencia de comandos; IPOPT y CasADi no se usan en esta implementación.

\[
x_{k+1}=x_k+v_k\cos(\theta_k)\Delta t,
\qquad
y_{k+1}=y_k+v_k\sin(\theta_k)\Delta t,
\]

\[
\theta_{k+1}=\theta_k+\omega_k\Delta t.
\]

La función de coste usada por el script penaliza error de distancia, error al
punto \(p^\star\), rumbo, esfuerzo, potencia de rueda y variación de comando:

\[
\begin{aligned}
J = \sum_{k=0}^{N-1}(&
q_d e_d(k)^2 + q_t\|p^\star_k-p_k\|^2
+ q_\psi e_\psi(k)^2 \\
&+ r_v v_k^2 + r_\omega \omega_k^2
+ r_{whl}(\omega_{L,k}^2+\omega_{R,k}^2)\\
&+ s_v\Delta v_k^2+s_\omega\Delta\omega_k^2)
+ q_f\|p^\star_N-p_N\|^2 .
\end{aligned}
\]

El coste incluye dos términos de posición con roles distintos: \(q_d e_d^2\)
penaliza el error escalar de distancia a Bill (mantiene el radio correcto
respecto al objetivo, pero no corrige la desviación lateral); \(q_t\|p^\star-p\|^2\)
penaliza la distancia euclídea al punto \(p^\star\), introduciendo una corrección
lateral suave que alinea el robot detrás de Bill. Ambos son necesarios: sin
\(q_t\) el robot mantiene el radio pero puede quedar desplazado lateralmente;
sin \(q_d\) el seguimiento del radio depende íntegramente de la posición de
\(p^\star\), con menor robustez ante errores de predicción de Bill.
El valor \(q_t=7{,}0\ll q_d=48{,}0\) asegura que la corrección lateral no
interfiera con la regulación de distancia.

Los pesos usados en el script fueron
\(q_d=48{,}0\), \(q_t=7{,}0\), \(q_\psi=1{,}2\),
\(r_v=1{,}4\), \(r_\omega=1{,}8\), \(r_{whl}=0{,}020\),
\(s_v=65{,}0\), \(s_\omega=18{,}0\) y \(q_f=6{,}0\), con horizonte
\(N=10\), \(\Delta t=\SI{0.22}{s}\) y dos iteraciones de mejora por ciclo.
Esta implementación reside en \texttt{pioneer\_pid\_follower\_controller.lua} y
aplica exclusivamente a la escena de seguimiento (\path{Actividad2_1_Basic_Follower.ttt}).
La escena warehouse (\path{Sim_T2_Phase1_Basic.ttt}) utiliza el controlador de tarea
\texttt{phase1\_r1\_task\_controller.lua}, que implementa los mismos cinco modos
(P, PI, PID, LQR y NMPC) con la misma ganancia $K$ y el mismo solucionador de
horizonte. La diferencia respecto al follower es que el objetivo en el warehouse
es un punto de navegación estático (no el punto $p^*$ detrás de Bill) y el
objetivo se reconstruye desde los errores de distancia y rumbo disponibles en el
planificador. La comparación de controladores presentada en este capítulo usa el
follower, donde el objetivo móvil ejercita mejor la respuesta dinámica.

La búsqueda local es determinista y de vecindario: la secuencia se inicializa
con los comandos del ciclo anterior desplazados un paso, y en cada pasada se
prueban perturbaciones \(\pm\Delta v\) y \(\pm\Delta\omega\) sobre cada elemento
del horizonte. En la primera pasada los pasos son \(\Delta v=0.065\) y
\(\Delta\omega=0.16\); en la segunda bajan a \(0.025\) y \(0.065\). Solo se
acepta una perturbación si reduce \(J\). Por eso el coste es predictivo y el
modelo es no lineal, pero el solucionador es una búsqueda local pequeña, no un
SQP ni un método interior-point.

Los límites de actuación se aplican durante la búsqueda:

\[
v_k\in[-v_{rev,max},v_{max}],
\qquad
\omega_k\in[-\omega_{max},\omega_{max}].
\]

Esta versión conserva el modelo no lineal y el coste predictivo, pero limita la
optimización a una búsqueda local pequeña sobre comandos. El alcance queda
limitado al método de disparo directo implementado en Lua.

\section{Resultados de seguimiento}

\begin{table}[H]
\centering
\footnotesize
\setlength{\tabcolsep}{3pt}
\begin{tabular}{@{}P{0.13\textwidth}P{0.16\textwidth}P{0.15\textwidth}P{0.15\textwidth}P{0.15\textwidth}P{0.12\textwidth}@{}}
\toprule
\textbf{Modo} & \textbf{Error mov. [m]} & \textbf{Error final [m]} & \textbf{Energía [J]} & \textbf{RMS \(\Delta u\)} & \textbf{Flips}\\
\midrule
P & 0,3542 & 0,0982 & \textbf{49,950} & 0,0475 & 11\\
PI & 0,2011 & 0,1018 & 64,419 & 0,1318 & 16\\
PID & 0,1915 & \textbf{0,0317} & 102,381 & 0,1759 & 26\\
LQR & 0,0231 & 0,0865 & 92,287 & 0,0873 & 12\\
NMPC & \textbf{0,0198} & 0,0887 & 77,174 & 0,0576 & \textbf{4}\\
\bottomrule
\multicolumn{6}{l}{\footnotesize \textit{n}=1 por modo. Energía: estimación cinemática aproximada (véase §\ref{sec:energy-model}). Indicadores cualitativos.}\\
\end{tabular}
\caption{Métricas del seguidor de Bill por modo de control.}
\label{tab:follower-controller-metrics}
\end{table}

\begin{figure}[H]
  \centering
  \includegraphics[width=\textwidth]{follower_pid_comparison.png}
  \caption{Comparación temporal de P, PI, PID, LQR y NMPC en la escena de seguimiento de Bill.}
  \label{fig:follower-pid-comparison}
\end{figure}

\begin{figure}[H]
  \centering
  \makebox[\textwidth][c]{\includegraphics[width=1.08\textwidth]{follower_wheel_power_comparison.png}}
  \caption{Coste físico aproximado del seguidor: potencia instantánea y energía acumulada en las ruedas.}
  \label{fig:follower-wheel-power}
\end{figure}

\clearpage

\begin{figure}[H]
  \centering
  \makebox[\textwidth][c]{\includegraphics[width=1.08\textwidth]{follower_wheel_command_smoothness.png}}
  \caption{Suavidad de comandos del seguidor: variación de las órdenes de rueda y cambios de sentido.}
  \label{fig:follower-wheel-smoothness}
\end{figure}

La Fig.~\ref{fig:follower-pid-comparison} muestra la respuesta temporal de los
cinco modos; la Fig.~\ref{fig:follower-wheel-power} compara potencia y energía
de rueda; la Fig.~\ref{fig:follower-wheel-smoothness} compara la suavidad de
comandos. P gasta menos energía y deja más error mientras Bill se mueve. PID
corrige bien al final pero paga con más variación de comando y más energía.
LQR y NMPC bajan el error medio durante el seguimiento; NMPC queda más suave
en inversiones de sentido y RMS de \(\Delta u\) porque incluye en el coste la
variación de comando y el esfuerzo de rueda.
